% Book-specific commands, environments and metadata.

% Title metadata used by the frontmatter pages.
\newcommand{\thebooktitle}{Từ RNN đến Transformer}
\newcommand{\thebooksubtitle}{Sáu bài báo gốc từ 1997 đến 2021, và những gì đến sau}
\newcommand{\thebookauthor}{Giang Dang}

% Inline code.
\newcommand{\code}[1]{\texttt{#1}}

% ---------------------------------------------------------------------------
% Glossed terms
% ---------------------------------------------------------------------------
% This book is written in Vietnamese for readers who will go on to read the
% papers in English. A term the book translates therefore carries its English
% original in parentheses at EVERY occurrence, not only the first: a reader who
% opens the book at chapter 9 gets the same help as one who started at page one.
%
%   \tn{mang hoi quy}{recurrent neural network}  ->  mang hoi quy (recurrent
%                                                    neural network)
%
% Terms the book keeps in English (attention, softmax, gradient, embedding,
% transformer) are typed plainly and take no parentheses, because "attention
% (attention)" helps nobody. Which terms fall on which side is settled once in
% appendix B and does not drift.
%
% Writing every gloss through this macro rather than by hand is what makes the
% rule checkable later: a term from the glossary appearing in the prose outside
% \tn is a miss, and that is a grep rather than a parser. See the SPEC's open
% items.
\newcommand{\tn}[2]{#1 (#2)}

% ---------------------------------------------------------------------------
% Mathematical notation
% ---------------------------------------------------------------------------
% One notation for the whole book (see chapter 1), with appendix A mapping it
% back to what each of the six papers wrote. These macros exist so that a
% decision about how a vector is set is made once here rather than 400 times in
% the chapters.

% Vectors are bold lowercase, matrices bold uppercase. Two macros with the same
% body, deliberately: the distinction is worth being able to grep for.
\newcommand{\vect}[1]{\mathbf{#1}}
\newcommand{\mat}[1]{\mathbf{#1}}

\newcommand{\transpose}{^{\mathsf{T}}}
\newcommand{\R}{\mathbb{R}}
\newcommand{\E}{\mathbb{E}}
\newcommand{\norm}[1]{\lVert #1 \rVert}
\newcommand{\abs}[1]{\lvert #1 \rvert}
\newcommand{\inner}[2]{\langle #1, #2 \rangle}

% Partial derivatives and Jacobians. Chapters 1 to 4 are mostly these.
\newcommand{\pd}[2]{\frac{\partial #1}{\partial #2}}
\newcommand{\jac}[2]{\frac{\partial #1}{\partial #2}}

\DeclareMathOperator{\softmax}{softmax}
\DeclareMathOperator{\attention}{Attention}
\DeclareMathOperator{\headop}{head}
\DeclareMathOperator{\multihead}{MultiHead}
\DeclareMathOperator{\layernorm}{LayerNorm}
\DeclareMathOperator{\ffn}{FFN}
\DeclareMathOperator{\diag}{diag}
\DeclareMathOperator*{\argmax}{arg\,max}
\DeclareMathOperator*{\argmin}{arg\,min}
\DeclareMathOperator{\tr}{tr}
\DeclareMathOperator{\rank}{rank}

\newcommand{\dmodel}{d_{\mathrm{model}}}
\newcommand{\dk}{d_{k}}
\newcommand{\dv}{d_{v}}
\newcommand{\dff}{d_{\mathrm{ff}}}

% ---------------------------------------------------------------------------
% Theorem environments
% ---------------------------------------------------------------------------
% One shared counter, numbered by chapter, so that a later chapter can cite
% "dinh ly 3.4" with no ambiguity about which sequence it belongs to. Babel's
% Vietnamese locale already supplies \proofname as "Chung minh".
\theoremstyle{definition}
\newtheorem{definition}{Định nghĩa}[chapter]
\newtheorem{example}[definition]{Ví dụ}

\theoremstyle{plain}
\newtheorem{theorem}[definition]{Định lý}
\newtheorem{lemma}[definition]{Bổ đề}
\newtheorem{proposition}[definition]{Mệnh đề}
\newtheorem{corollary}[definition]{Hệ quả}

\theoremstyle{remark}
\newtheorem{remark}[definition]{Nhận xét}

% The algorithm float is the one caption label in this book that babel does not
% localize, because `algorithm` sets its own name rather than taking one from
% the language. Left alone it prints "Algorithm 1" in the middle of a Vietnamese
% page. Chapter 3 shipped that way; chapter 5 would have been the second.
\floatname{algorithm}{Thuật toán}

% ---------------------------------------------------------------------------
% The two recurring boxes
% ---------------------------------------------------------------------------
% Grey and rules only, no colour: these books are likely to be read in mono.

% Bridge box. The six papers do not form a chain on their own - Bahdanau builds
% on Cho rather than on Sutskever, the forget gate arrives three years after the
% 1997 LSTM, ViT stands on BERT and GPT. A bridge box is where an intermediate
% paper gets its half page: what it changed, and why the next chapter would not
% make sense without it. The argument of the chapter runs outside these boxes,
% so a reader can skip every one of them and still follow the spine.
\newtcolorbox{bridge}[1]{
  breakable, enhanced, sharp corners,
  colback = black!3, colframe = black!55,
  boxrule = 0.4pt, left = 3mm, right = 3mm, top = 2mm, bottom = 2mm,
  title = {Cầu nối: #1},
  fonttitle = \bfseries\small,
  coltitle = black, colbacktitle = black!12,
}

% Measured note. This book prints two kinds of number: what a paper reported,
% which is cited, and what I ran myself, which lives in one of these. The box
% exists so the difference is visible on the page rather than buried in a
% subordinate clause. Every number inside one traces to a file under research/.
\newtcolorbox{measured}[1]{
  breakable, enhanced, sharp corners,
  colback = white, colframe = black!70,
  boxrule = 0.8pt, left = 3mm, right = 3mm, top = 2mm, bottom = 2mm,
  title = {Tôi đo được: #1},
  fonttitle = \bfseries\small,
  coltitle = black, colbacktitle = black!8,
}

% ---------------------------------------------------------------------------
% End-of-chapter exercises
% ---------------------------------------------------------------------------
% Three tiers, one per goal the book set itself: understand the idea, apply it
% to code, push past it. Starred headings throughout - nothing cross-references
% an exercise, and the tiers would otherwise flood the table of contents with
% forty-eight entries.
\newcommand{\exercises}{\section*{Bài tập cuối chương}}
\newcommand{\tierunderstand}{\subsection*{Kiểm tra hiểu}}
\newcommand{\tierapply}{\subsection*{Để tay vào}}
\newcommand{\tierextend}{\subsection*{Đẩy xa hơn}}

\newlist{bt}{enumerate}{1}
\setlist[bt]{label = \arabic*., leftmargin = *, itemsep = 0.4em, topsep = 0.4em}
